\section{Introduction}
\label{sec:introduction}

\subsection{The Gap Between Structure-Layer and Search-Layer VRA Enforcement}

Redistricting litigation under Section~2 of the Voting Rights
Act~\citep{gingles1986} requires two distinct capabilities: (1) generating a
plan that satisfies VRA---that is, one that preserves majority-minority
districts in states where they are mandated---and (2) producing an ensemble
of valid alternative plans to establish what the space of compliant plans looks
like.
The first capability is addressed at the \emph{structure layer}: \vraSection{}
(D.0~\citep{dellaLibera2026vraSection}) biases the METIS bisection to produce
a single VRA-compliant plan.
The second capability requires a different tool: a Markov chain that samples the
\emph{conditional distribution} over plans that are simultaneously valid (equal
population, contiguity) and VRA-compliant.

Standard ensemble methods do not provide this.
\recom{}~\citep{deford2021}, Forest \recom{}~\citep{mccartan2020,dellaLibera2026forestRecom},
and Merge-Split all sample the full space of valid plans without distinguishing
VRA-compliant from non-compliant alternatives.
When used for litigation, these chains produce ensembles that include plans
with no majority-minority districts at all---plans that could never be legally
enacted in VRA-covered states.
Drawing conclusions about VRA compliance from such an ensemble conflates the
space of valid plans with the space of legally permissible plans.

This paper introduces \vr{}: a search-layer algorithm that samples the
conditional distribution over VRA-compliant redistricting plans.
The core idea is simple: extend \fr{} with a hard rejection rule that
unconditionally rejects any proposed step that would reduce the minority VAP
fraction of a protected district below the VRA threshold.
VRA compliance is enforced as an invariant of the chain, not as a soft
preference in the Metropolis-Hastings ratio.

\subsection{Two Contributions}

\begin{enumerate}
  \item \textbf{Hard-rejection VRA invariant.}
        \vr{} enforces VRA compliance as a chain invariant: every plan in
        the ensemble is guaranteed to preserve all protected districts above
        \texttt{vap\_threshold}.
        The protection set is fixed at chain construction from the initial
        plan, matching the legal standard that the submitted plan's
        majority-minority districts must be preserved in the alternative
        ensemble.
        (Section~\ref{sec:algorithm}.)

  \item \textbf{Three-way step accounting for audit.}
        Every step is categorised as accepted, VRA-rejected, or MH-rejected,
        with the identity $\text{accepted} + \text{vra\_rejections} +
        \text{mh\_rejections} = \text{steps\_taken}$ holding exactly.
        The audit chain records \texttt{protected\_districts},
        \texttt{vra\_rejection\_rate}, and the seed formula so that any
        verifier can replay the chain and confirm every VRA rejection.
        (Section~\ref{sec:audit}.)
\end{enumerate}

\subsection{Relationship to VRASection}

\vraSection{} is a structure-layer algorithm: it modifies METIS edge weights
to bias the initial bisection toward VRA-compliant splits.
\vr{} is a search-layer algorithm: it samples the neighbourhood of an
already-compliant plan.
They are complementary.
The recommended workflow is:
\begin{enumerate}
  \item Run \vraSection{} (\texttt{--structure ratio-optimal-vra}) to generate a
        VRA-compliant starting plan.
  \item Run \vr{} (\texttt{--search vra-recom}) to sample the distribution over
        compliant alternatives, with the \vraSection{} plan as the initial state.
\end{enumerate}
This ensures both the starting plan and the entire ensemble are VRA-compliant,
satisfying the legal requirement that the submitted plan and all proposed
alternatives respect Section~2.

\subsection{Paper Structure}

Section~\ref{sec:background} reviews the Gingles factors, the legal definition
of majority-minority districts, and the structural distinction between
structure-layer and search-layer VRA enforcement.
Section~\ref{sec:algorithm} gives the formal \vr{} algorithm, including the
three-way step accounting and the seed specification.
Section~\ref{sec:empirical} presents empirical results on North Carolina,
Georgia, and Texas, including VRA rejection rates and EC distribution
comparisons.
Section~\ref{sec:audit} describes the audit chain and its legal properties.
Section~\ref{sec:conclusion} concludes with usage guidance and open questions.
